\chapter{Μεθοδολογία και πειραματική διάταξη}
\label{ch:methodology}

\section{Σύστημα αναφοράς}
\label{sec:meth-machine}

Όλες οι μετρήσεις της εργασίας προέρχονται από ένα και μόνο σύστημα, το οποίο
φέρει και τις τρεις αρχιτεκτονικές που συγκρίνονται: έναν επεξεργαστή γενικού
σκοπού (\tl{CPU}), μια διακριτή μονάδα γραφικών (\tl{GPU}) και μια μονάδα νευρωνικής 
επεξεργασίας (\tl{NPU}) ενσωματωμένη στο ίδιο πακέτο με τον επεξεργαστή. Η επιλογή 
αυτή απομονώνει τις διαφορές των υλοποιήσεων από τις διαφορές των πλατφορμών, 
εισγάγει όμως έναν περιορισμό που καθορίζει τις συνθήκες μέτρησης: το σύστημα 
είναι φορητό και οι τρεις μονάδες μοιράζονται ένα κοινό θερμικό και ενεργειακό όριο. 
Ο χειρισμός του περιορισμού αυτού περιγράφεται στην ενότητα~\ref{sec:meth-conditions}.

\begin{table}[H]
    \centering
    \begin{tabular}{ll}
        \hline
        Σύστημα & \tl{Lenovo Legion 5 15AHP10} \\
        Λειτουργικό σύστημα & \tl{Arch Linux x86\_64}, πυρήνας \tl{7.1.5-arch1-2} \\
        Κύρια μνήμη & \qty{32}{\gibi\byte} \tl{DDR5-5600}, δύο κανάλια \\
        \hline
        \tl{CPU} & \tl{AMD Ryzen 7 260} \\
        Πυρήνες / νήματα & 8 / 16, έως \qty{5.10}{\giga\hertz} \\
        Κρυφή μνήμη & \qty{8}{\mebi\byte} \tl{L2}, \qty{16}{\mebi\byte} \tl{L3} \\
        Διανυσματικές επεκτάσεις & \tl{AVX-512} \\
        \hline
        \tl{GPU} & \tl{NVIDIA GeForce RTX 5060 Laptop} \\
        Μνήμη & \qty{8}{\gibi\byte} \\
        Λογισμικό & \tl{driver} \tl{610.43.03}, \tl{CUDA 13.3} \\
        Όρια ισχύος & \qty{50}{\watt} προεπιλογή, \qty{115}{\watt} μέγιστο \\
        \hline
        \tl{NPU} & \tl{AMD XDNA} (\tl{RyzenAI-npu1}) \\
        Λογισμικό & \tl{XRT 2.21.75}, \tl{firmware 1.5.5.391} \\
        \hline
    \end{tabular}
    \caption{Το σύστημα αναφοράς στο οποίο εκτελέστηκαν όλες οι μετρήσεις.}
    \label{tab:machine}
\end{table}

Ο επεξεργαστής διαθέτει οκτώ φυσικούς πυρήνες με ταυτόχρονη πολυνηματικότητα.
Όλες οι μετρήσεις εκτελούνται με οκτώ νήματα, ένα ανά φυσικό πυρήνα, ώστε να μη
μοιράζονται δύο νήματα την ίδια διανυσματική μονάδα. Η ιεραρχία κρυφής μνήμης
καθορίζει τα σημεία καμπής της μετρημένης κλίμακας εύρους ζώνης
(ενότητα~\ref{sec:meth-roofline}): ανάλογα με το μέγεθος της εισόδου και τον
τύπο δεδομένων, το σύνολο εργασίας μπορεί να είναι κατοικήσιμο στην \tl{L2},
στην \tl{L3} ή μόνο στην κύρια μνήμη, και η επίδοση κάθε υλοποίησης συγκρίνεται
με το ταβάνι που αντιστοιχεί στο δικό της σύνολο εργασίας.

Η κύρια μνήμη είναι διατεταγμένη σε δύο κανάλια, με δύο πανομοιότυπης
χωρητικότητας μονάδες των \qty{16}{\gibi\byte} σε ρυθμό \tl{5600 MT/s}. Η
λεπτομέρεια αυτή είναι σημαντική, καθώς η διάταξη ενός μόνο καναλιού υποδιπλασιάζει
το διαθέσιμο εύρος ζώνης, οπότε κάθε υλοποίηση που περιορίζεται από τη μνήμη θα
έδινε ουσιωδώς διαφορετικά αποτελέσματα σε ένα κατά τα άλλα ίδιο σύστημα.

Η μονάδα νευρωνικής επεξεργασίας δεν διαθέτει δική της μνήμη· τροφοδοτείται από
την ίδια \tl{DDR5} και μέσω των ίδιων ελεγκτών μνήμης με τον επεξεργαστή, στο
ίδιο πακέτο πυριτίου. Αυτό έχει δύο συνέπειες που επανεμφανίζονται σε όλη την
εργασία: το κόστος προετοιμασίας των δεδομένων στον \tl{host} δεν είναι
διαχωρίσιμο από το κόστος του υπολογισμού, και ο μετρητής ενέργειας που
χρησιμοποιείται καλύπτει ολόκληρο το πακέτο και όχι τα πλακίδια \tl{AIE}
ξεχωριστά (ενότητα~\ref{sec:meth-energy}).

Τα στοιχεία του πίνακα~\ref{tab:machine} έχουν επαληθευτεί από το ίδιο το
σύστημα και όχι από φυλλάδια κατασκευαστή. Χαρακτηριστικά που δεν είναι άμεσα
μετρήσιμα, όπως η ονομασία της μικροαρχιτεκτονικής, το θεωρητικό εύρος ζώνης ή
οι διαφημιζόμενες επιδόσεις σε \tl{TOPS}, δεν χρησιμοποιούνται πουθενά στη
σύγκριση. Κάθε ταβάνι εύρους ζώνης και υπολογιστικής ικανότητας με το οποίο
συγκρίνονται οι υλοποιήσεις μετριέται στο ίδιο μηχάνημα, με τη διαδικασία της
ενότητας~\ref{sec:meth-roofline}.

\section{Δομή της υποδομής μετρήσεων}
\label{sec:meth-harness}

Η υποδομή μετρήσεων οργανώνεται σε τέσσερα επίπεδα, με μία αρχή να διατρέχει
τον σχεδιασμό της: η διαδικασία της μέτρησης είναι κοινή για όλες τις
αρχιτεκτονικές και μόνο ο υπολογιστικός πυρήνας διαφέρει. Καμία υλοποίηση δεν
ελέγχει τον τρόπο με τον οποίο μετριέται η ίδια, ώστε να μην μπορεί καμία να
χρονομετρεί, ακούσια, κάτι διαφορετικό από τις υπόλοιπες.

Το πρώτο επίπεδο είναι ένας κοινός πυρήνας σε \tl{C++} τον οποίο μοιράζονται και
οι τρεις υλοποιήσεις. Σε αυτόν ανήκουν η ανάγνωση των παραμέτρων, η παραγωγή της
εισόδου, το χρονοδιάγραμμα των επιπέδων του διτονικού δικτύου, η ανάγνωση των
μετρητών ενέργειας, ο βρόχος των επαναλήψεων προθέρμανσης και μέτρησης, ο
έλεγχος ορθότητας και η εκτύπωση της αναφοράς. Επειδή ο βρόχος αυτός υπάρχει σε
ένα μόνο σημείο, οι χρονομετρήσεις και οι μετρήσεις ενέργειας όλων των
αρχιτεκτονικών ορίζονται από τον ίδιο κώδικα.

Το δεύτερο επίπεδο είναι τα εκτελέσιμα ανά αρχιτεκτονική. Κάθε ένα υλοποιεί την
ίδια διεπαφή για τους τρεις αλγορίθμους που συγκρίνονται, δηλαδή τη διτονική
ταξινόμηση, το \tl{map-reduce} με τοπικούς σωρούς και την υλοποίηση αναφοράς,
παρέχοντας αποκλειστικά τη συνάρτηση που εκτελεί τον υπολογισμό. Η συνάρτηση
αυτή καλείται από τον κοινό βρόχο μέτρησης του πρώτου επιπέδου.

Το τρίτο επίπεδο είναι ο οδηγός των σαρώσεων. Απαριθμεί τον παραμετρικό χώρο της
ενότητας~\ref{sec:meth-parameters} και εκτελεί κάθε περίπτωση ως ξεχωριστή
διεργασία, καταγράφοντας αυτούσια την έξοδό της μαζί με την έκβασή της. Η
εκτέλεση σε ξεχωριστές διεργασίες εξασφαλίζει ότι κάθε περίπτωση ξεκινά από
καθαρή κατάσταση μνήμης και ότι η αποτυχία μιας δεν ακυρώνει τη σάρωση.

Το τέταρτο επίπεδο μετατρέπει τα ακατέργαστα αποτελέσματα σε πίνακες και
διαγράμματα. Η ανάλυση γίνεται πάνω στο ίδιο κείμενο αναφοράς που διαβάζει και ο
χρήστης: δεν υπάρχει χωριστή διαδρομή εξόδου προς τα εργαλεία, οπότε οι τιμές που
εμφανίζονται στα διαγράμματα είναι εξ ορισμού οι τιμές που τύπωσε το εκτελέσιμο.
Μαζί με τα αποτελέσματα αποθηκεύεται και το περιβάλλον εκτέλεσης, όπως
περιγράφεται στην ενότητα~\ref{sec:meth-conditions}. Οι εντολές που παράγουν
κάθε στάδιο δίνονται στο παράρτημα~\ref{app:reproducibility}.

\section{Μέτρηση χρόνου}
\label{sec:meth-timing}

\section{Μέτρηση ενέργειας}
\label{sec:meth-energy}

\subsection{Πεδία μέτρησης: \tl{algo} και \tl{end-to-end}}
\label{sec:meth-energy-scope}

\section{Μοντέλο \tl{roofline} και μέτρηση κίνησης δεδομένων}
\label{sec:meth-roofline}

\section{Υλοποιήσεις αναφοράς}
\label{sec:meth-baselines}

\subsection{Αναφορά στον \tl{host}}
\label{sec:meth-baseline-host}

\subsection{Βιβλιοθήκες σύγκρισης}
\label{sec:meth-baseline-libraries}

\section{Έλεγχος ορθότητας}
\label{sec:meth-correctness}

\section{Παραμετρικός χώρος}
\label{sec:meth-parameters}

\subsection{Τύποι δεδομένων}
\label{sec:meth-dtypes}

\subsection{Κατανομές εισόδου}
\label{sec:meth-distributions}

\section{Συνθήκες μέτρησης και επαναληψιμότητα}
\label{sec:meth-conditions}

Το σύστημα δεν βρίσκεται εκ των προτέρων σε κατάσταση κατάλληλη για μέτρηση. Στην
προεπιλεγμένη διαμόρφωσή του ο διαχειριστής συχνότητας (\tl{governor}) είναι ο
\tl{powersave} και η ενίσχυση συχνότητας (\tl{boost}) ενεργή, οπότε η συχνότητα
λειτουργίας μεταβάλλεται ανάλογα με το φορτίο και τη θερμοκρασία. Η μεταβολή
αυτή δεν είναι τυχαίος θόρυβος γύρω από μια σταθερή τιμή, είναι συστηματική
μετατόπιση κατά τη διάρκεια της ίδιας σάρωσης, καθώς το σύστημα θερμαίνεται.
Επομένως η τυπική απόκλιση ενός συνόλου δειγμάτων που ελήφθησαν ενώ οι συνθήκες
μετατοπίζονταν δεν αποτυπώνει την αβεβαιότητα της μέτρησης, και η αναφορά μέσης
τιμής και απόκλισης δεν αρκεί από μόνη της για να θεωρηθεί η μέτρηση αξιόπιστη.

Για τον λόγο αυτό, πριν από κάθε σάρωση το σύστημα καθηλώνεται σε μια σταθερή
και τεκμηριωμένη κατάσταση, η οποία συνοψίζεται στον πίνακα~\ref{tab:conditions}.
Η κατάσταση αυτή καταγράφεται πριν από τη σάρωση και αποκαθίσταται μετά το πέρας
της. Οι αντίστοιχες εντολές δίνονται στο παράρτημα~\ref{app:reproducibility}.

\begin{table}[H]
    \centering
    \begin{tabular}{ll}
        \hline
        Παράμετρος & Κατάσταση κατά τη μέτρηση \\
        \hline
        Διαχειριστής συχνότητας \tl{CPU} & \tl{performance} \\
        Ενίσχυση συχνότητας \tl{CPU} & απενεργοποιημένη \\
        Νήματα \tl{CPU} & 8, ένα ανά φυσικό πυρήνα \\
        \tl{GPU persistence mode} & ενεργοποιημένο \\
        Όριο ισχύος \tl{GPU} & σταθερό σε όλη τη σάρωση \\
        Περίοδος ψύξης μεταξύ μονάδων & \qty{15}{\second} \\
        Επαναλήψεις προθέρμανσης & 5 ανά περίπτωση \\
        Μετρούμενες επαναλήψεις & 50 ανά περίπτωση \\
        \hline
    \end{tabular}
    \caption{Οι συνθήκες στις οποίες καθηλώνεται το σύστημα πριν από κάθε σάρωση.}
    \label{tab:conditions}
\end{table}

Το όριο ισχύος της μονάδας γραφικών επιλέγεται μία φορά και διατηρείται σταθερό
καθ'όλη τη μέτρηση της εργασίας. Η συγκεκριμένη τιμή είναι κατά βάση αυθαίρετη,
εφόσον βρίσκεται εντός των ορίων του πίνακα~\ref{tab:machine}. Εκείνο που έχει
σημασία είναι ότι δεν μεταβάλλεται ούτε μεταξύ των υλοποιήσεων ούτε κατά τη
διάρκεια της σάρωσης, καθώς κάθε αλλαγή της θα μετέβαλλε ταυτόχρονα τον χρόνο
εκτέλεσης και την ενέργεια που καταναλώνει η μονάδα.

Όπως αναφέρθηκε στην ενότητα~\ref{sec:meth-machine}, οι τρεις μονάδες μοιράζονται
ένα κοινό θερμικό όριο. Μεταξύ διαδοχικών μονάδων μεσολαβεί περίοδος ψύξης
\qty{15}{\second}, ώστε η κάθε μια να μην ξεκινά τη μέτρησή της από τη θερμική
κατάσταση που άφησε η προηγούμενη, ενώ κατά τη διάρκεια μιας σάρωσης το σύστημα
δεν εκτελεί άλλο υπολογιστικό φορτίο. Πρόκειται για μετριασμό και όχι για
εξάλειψη του προβλήματος: σε ένα φορητό σύστημα η θερμική σύζευξη μεταξύ των
μονάδων δεν μπορεί να αρθεί και αυτό αποτελεί αναγνωρισμένο περιορισμό των
μετρήσεων.

Κάθε περίπτωση εκτελείται 5 φορές ως προθέρμανση και στη συνέχεια 50 φορές υπό
μέτρηση. Αναφέρονται η μέση τιμή και η τυπική απόκλιση των 50 δειγμάτων, καθώς
και το ελάχιστο ως ένδειξη της βέλτιστης περίπτωσης. Πριν από κάθε επανάληψη η
είσοδος ανακτάται από ένα αντίγραφο ασφαλείας, εκτός της χρονομετρούμενης
περιοχής. Αυτό εξυπηρετεί δύο σκοπούς: το κόστος της αντιγραφής δεν προσμετράται
στον μετρούμενο χρόνο, και κάθε επανάληψη ξεκινά από πανομοιότυπη είσοδο αντί να
επεξεργάζεται τη μερικώς διατεταγμένη έξοδο της προηγούμενης. Η δεύτερη συνθήκη
είναι κρίσιμη, διότι οι υλοποιήσεις που εξαρτώνται από τα δεδομένα θα
εμφάνιζαν τεχνητά βελτιωμένους χρόνους σε μια ήδη μερικώς ταξινομημένη είσοδο.

Πέρα από τις επαναλήψεις εντός μιας περίπτωσης, κάθε συνδυασμός κατανομής
εισόδου και σπόρου (\tl{seed}) γεννήτριας τυχαίων αριθμών αποτελεί ανεξάρτητο δείγμα. Τα
διαγράμματα της εργασίας συναθροίζουν αυτά τα δείγματα και απεικονίζουν τη
διασπορά τους ως ζώνες σφάλματος, ώστε η αβεβαιότητα που οφείλεται στα δεδομένα
να διακρίνεται από την αβεβαιότητα που οφείλεται στη μέτρηση.

Τέλος, κάθε εκτέλεση καταγράφει η ίδια το περιβάλλον στο οποίο έτρεξε: την
αναθεώρηση του αποθετηρίου, τον πυρήνα, τον μεταγλωττιστή και τις εκδόσεις των
οδηγών και των βιβλιοθηκών εκτέλεσης κάθε μονάδας. Οι πληροφορίες αυτές
αποθηκεύονται μαζί με τα αποτελέσματα, ώστε κάθε σύνολο μετρήσεων να
τεκμηριώνει από μόνο του τις συνθήκες παραγωγής του και να μην εξαρτάται από
εξωτερική καταγραφή που ενδέχεται να αποκλίνει.
